%!TEX program = xelatex
\documentclass[aspectratio=169]{beamer}
\usepackage[backend=biber,style=numeric-comp,sorting=none]{biblatex}
\addbibresource{biblio.bib}

\usecolortheme[light,accent=orange]{solarized}

\usepackage{subcaption}
\usepackage{gitdags}

\usepackage{blindtext}
\usepackage{xltxtra}
\usepackage{wrapfig}

\usepackage{listings}
\usepackage{multicol}
\usepackage{verbatim}

\usepackage{graphicx,changepage}

\newcommand{\adjustimg}{% Horizontal adjustment of image
  \checkoddpage%
  \ifoddpage\hspace*{\dimexpr\evensidemargin-\oddsidemargin}\else\hspace*{-\dimexpr\evensidemargin-\oddsidemargin}\fi%
}
\newcommand{\centerimg}[2][width=\textwidth]{% Center an image
  \makebox[\textwidth]{\adjustimg\includegraphics[#1]{#2}}%
}

\defaultfontfeatures{Ligatures=TeX}

\title{Oh my GIT!}
\subtitle{Git ate my work and other stories}
\author{Vojtěch Vladyka}
\date{\today}

% \setbeamercolor
%
% solarizedBase03
% solarizedBase02
% solarizedBase01
% solarizedBase00
% solarizedBase0
% solarizedBase1
% solarizedBase2
% solarizedBase3
% solarizedYellow
% solarizedOrange
% solarizedRed
% solarizedMagenta
% solarizedViolet
% solarizedBlue
% solarizedCyan
% solarizedGreen

\begin{document}
    \frame{\titlepage}
    \begin{frame}
       \frametitle{Table of contents}
       \begin{enumerate}
           \item What is git
           \\   \textcolor{solarizedRebase01}{\footnotesize\hspace{1em} ...and why you should care}	
           \item Setup git
           \\   \textcolor{solarizedRebase01}{\footnotesize\hspace{1em} Clients, differences, common issues}	
           \item Git essentials
           \\   \textcolor{solarizedRebase01}{\footnotesize\hspace{1em} Git controls crash course}	
           \item Scenarios
           \\   \textcolor{solarizedRebase01}{\footnotesize\hspace{1em} We are going to break things there. A lot.}	
           \item Q \& A
           \\   \textcolor{solarizedRebase01}{\footnotesize\hspace{1em} Umm, why was I here again?}
       \end{enumerate}
    \end{frame}

    \begin{frame}
        \frametitle{What is git?}
        \begin{itemize}
            \item Git is distributed by design
            \item Works with whole sourcetree (unlike SVN which works with individual files and their revisions) and always capture state of working tree
            \item Strong support for non-linear development $\rightarrow$ rapid branching \& merging
            \item Cryptographic authentication of history - every commit has SHA-1 hash of everything leading to that point
            \item Version control system made by Linus Torvalds (mostly) at 2005 for Linux Kernel development
            \item Supports rapid branching \& merging
            \item it is pronounced GIT with G like GIF ~\footfullcite{linus-speech}
        \end{itemize}
    \end{frame}

    \begin{frame}[fragile]
        \frametitle{Git philosophy}
        \begin{multicols}{2}
            \begin{itemize}
                \item Create branch for each feature
                \item Commit often
                \item Merge back to master as soon as possible
                \item Create small isolated as-much-as-possible pull requests
            \end{itemize}
            \includegraphics[width=0.45\textwidth]{pull-request.png}
        \end{multicols}
    \end{frame}
    
    \begin{frame}
        \frametitle{Setup git}
        \framesubtitle{Git clients}
        \begin{enumerate}
                % icons for platforms
                % add short desc to each
            \item Git (all platforms)
            \item Git Extensions (Windows, Mac)
            \item Tortoise Git (Windows)
            \item Sourcetree (Windows)
            \item Fork (Windows, Mac)
            \item Gitkraken (all platforms)
            \item ...many others
        \end{enumerate}
    \end{frame}
    \begin{frame}
        \frametitle{Setup git}
        \framesubtitle{Git Extensions (Git client)}
        \noindent\centerimg[width=0.8\paperwidth]{gitex.png} 
        Download it here $\rightarrow$ \url{https://gitextensions.github.io}
    \end{frame}

    \begin{comment}
    \begin{frame}
        \begin{tikzpicture}
            % Commit DAG
            \gitDAG[grow right sep = 2em]{
                A -- B -- {
                  C,
                  D -- E,
                }
            };
            % Tag reference
            \gittag
                [v0p1]      % node name
                {v0.1}      % node text
                {above=of A}% node placement
                {A}         % target
            % Remote branch
            \gitremotebranch
                [originmaster]
                {origin/master}
                {above=of C}
                {C}
            % Branch
            \gitbranch
                {master}
                {above=of E}
                {E}
            % HEAD reference
            \gitHEAD
                {above=of master}
                {master}
        \end{tikzpicture}
    \end{frame}
    \end{comment}
    \begin{frame}[fragile]
        \frametitle{Git essentials}
        \framesubtitle{Short and incomplete command overview}
        \begin{multicols}{2}
        \begin{lstlisting}[gobble=12]
            git init                
                clone

                fetch
                merge
                push
                pull

                branch
                checkout
                merge
                branch -d


            git log                
                log --follow [file]
                reflog

                add
                commit

                reset
                reset --hard
                stash

                rebase
                cherrypick
        \end{lstlisting}
        \end{multicols}
    \end{frame}
    \begin{frame}
        \frametitle{Git essentials}
        \noindent\centerimg[width=0.4\textwidth]{git.png}
    \end{frame}
    \begin{frame}
        \frametitle{Git essentials}
        \framesubtitle{Workflow overview}
        \noindent\centerimg[width=0.8\textwidth]{workflow.png}
    \end{frame}
    \begin{frame}[fragile]
        \frametitle{Git essentials}
        \framesubtitle{Init repo}
        \begin{lstlisting}[gobble=12]
            git init
        \end{lstlisting}
        $\rightarrow$ initialize new repository in current directory with all its contents
        \begin{lstlisting}[gobble=12]
            git clone https://somerepo.com/repo.git
        \end{lstlisting}
        $\rightarrow$ downloads repository with all current branches \& commits
    \end{frame}
    \begin{frame}[fragile]
        \frametitle{Git essentials}
        \framesubtitle{Fetching changes and returning them back}
        \begin{lstlisting}[gobble=12]
            git pull
        \end{lstlisting}
        $\rightarrow$ Download commits from server and apply them to your working tree
        \begin{lstlisting}[gobble=12]
            git push
        \end{lstlisting}
        $\rightarrow$ Upload your commits to server
        \begin{lstlisting}[gobble=12]
            git fetch --all
        \end{lstlisting}
        $\rightarrow$ Download commits from server (from all branches eventually)
        \begin{lstlisting}[gobble=12]
            git commit [files] -m "Commit message"
        \end{lstlisting}
        $\rightarrow$ Create a commit with commit message and specified files\newline
        $\rightarrow$ You can use -a to commit all changes instead list of files
        \begin{lstlisting}[gobble=12]
            git add [file]
        \end{lstlisting}
        $\rightarrow$ Move file to staging area - add it to commit list
    \end{frame}
    \begin{frame}
        \frametitle{Git essentials}
        \noindent\centerimg[width=0.6\textwidth]{git_commit.png}
    \end{frame}
    \begin{frame}[fragile]
        \frametitle{Git essentials}
        \framesubtitle{Undoing changes}
        \begin{lstlisting}[gobble=12]
            git reset
        \end{lstlisting}
        $\rightarrow$ Soft reset - remove files from staging area and move them back to working area. Nothing is actually reverted, this resets it only for git.
        \begin{lstlisting}[gobble=12]
            git reset --hard
        \end{lstlisting}
        $\rightarrow$ Hard reset - this do all what soft reset but also REMOVES your uncommitted changes.
        \newline\newline
        $\rightarrow$ You can do this only locally. Once you push your changes, you cannot get it back. (not exactly true but you will break a lot of thing for all colleagues). Only way how to propperly revert things it using git revert.
    \end{frame}

    \begin{frame}
        \noindent\centerimg[height=1.0\textheight]{it-hell.png}
    \end{frame}

    \begin{frame}[fragile]
        \frametitle{Git essentials}
        \framesubtitle{Branching}
        \begin{lstlisting}[gobble=12]
            git checkout -b feature1
            (git branch feature1 ; git checkout feature1)
        \end{lstlisting}
        \resizebox{!}{0.25\paperheight}{%
        \begin{tikzpicture}
            % Commit DAG
            \gitDAG[grow right sep = 2em]{
                A -- B -- C -- {
                    D,
                    E -- F
                }
            };
            % Branch
            \gitbranch
                {master}
                {above=of D}
                {D}
            % Branch
            \gitbranch
                {feature1}
                {above=of F}
                {F}
            % HEAD reference
            \gitHEAD
                {above=of feature1}
                {feature1}
        \end{tikzpicture}%
        }
        \begin{lstlisting}[gobble=12]
            git checkout master
            git merge feature1
        \end{lstlisting}
        \resizebox{!}{0.25\paperheight}{%
        \begin{tikzpicture}
            % Commit DAG
            \gitDAG[grow right sep = 2em]{
                A -- B -- C -- {
                    D,
                    E -- F
                } -- G
            };
            % Branch
            \gitbranch
                {master}
                {below=of G}
                {G}
            % Branch
            \gitbranch
                {feature1}
                {below=of F}
                {F}
            % HEAD reference
            \gitHEAD
                {below=of master}
                {master}
        \end{tikzpicture}%
        }

    \end{frame}
    \begin{frame}
        \noindent\centerimg[width=1.0\textwidth]{merging-branches-2.png}
    \end{frame}
    \begin{frame}[fragile]
        \frametitle{Git essentials}
        \framesubtitle{Logging and searching lost stuff}
        \begin{lstlisting}[gobble=12]
            git log --graph
        \end{lstlisting}
        $\rightarrow$ Shows log for whole repository tree
        \begin{lstlisting}[gobble=12]
            git log --follow [file]
        \end{lstlisting}
        $\rightarrow$ Shows log for selected file even across renames
        \begin{lstlisting}[gobble=12]
            git reflog
        \end{lstlisting}
        $\rightarrow$ Shows log of ALL changes in repository. Stashes, detached branches, all...
    \end{frame}
    \begin{frame}[fragile]
        \frametitle{Git essentials}
        \framesubtitle{Stashing - let me just put this here for later (and never use it again)}
        \begin{lstlisting}[gobble=12]
            git stash
            git stash push
        \end{lstlisting}
        $\rightarrow$ Stashes current changes to new stash stack
        \begin{lstlisting}[gobble=12]
            git stash pop
        \end{lstlisting}
        $\rightarrow$ Pops newest stash and apply it to current working copy
        \begin{lstlisting}[gobble=12]
            git stash list
        \end{lstlisting}
        $\rightarrow$ Show all stashes in list
        \begin{lstlisting}[gobble=12]
            git stash apply stash@{2}
        \end{lstlisting}
        $\rightarrow$ Get stash \#2 and apply it to current working copy
        \begin{lstlisting}[gobble=12]
            git checkout stash@{2} -- somefile
        \end{lstlisting}
        $\rightarrow$ Get stash \#2 and apply it to specified files in current working copy
    \end{frame}
    \begin{frame}[fragile]
        \frametitle{Git essentials}
        \framesubtitle{Submodules}
        \begin{lstlisting}[gobble=12]
            git submodule add <url to submodule repository>
        \end{lstlisting}
        $\rightarrow$ This adds new submodule to current repository and also clones it.\newline
        \begin{lstlisting}[gobble=12]
            git submodule init
        \end{lstlisting}
        $\rightarrow$ This clones submodules when you just cloned new repository with submodules.\newline
        \begin{lstlisting}[gobble=12]
            git submodule update --init --recursive
        \end{lstlisting}
        $\rightarrow$ And this downloads submodules to version specified in repository (and eventually clones new submodules)\newline
        \begin{lstlisting}[gobble=12]
            git submodule foreach 'do something'
        \end{lstlisting}
        $\rightarrow$ This runs 'do something' on every submodule.\newline
    \end{frame}
    \begin{frame}[fragile]
        \frametitle{Git essentials}
        \framesubtitle{Rebase}
        \begin{lstlisting}[gobble=12]
            git rebase target_branch source_branch
        \end{lstlisting}
        $\rightarrow$ Take source branch and reattach it on end of target branch. This creates series of new commits and remove old ones from tree (althrough they will still be accessible by reflog)\newline
        \resizebox{!}{0.25\paperheight}{%
        \begin{tikzpicture}
            % Commit DAG
            \gitDAG[grow right sep = 2em]{
                A -- B --{
                    C -- D ,
                    E -- F -- G
                }
            };
            % Branch
            \gitbranch
                {master}
                {above=of D}
                {D}
            % Branch
            \gitbranch
                {feature1}
                {above=of G}
                {G}
            % HEAD reference
            \gitHEAD
                {above=of feature1}
                {feature1}
        \end{tikzpicture}%
        }
        \begin{center}$\Downarrow$\end{center}
        \resizebox{!}{0.25\paperheight}{%
        \begin{tikzpicture}
            % Commit DAG
            \gitDAG[grow right sep = 2em]{
                A -- B -- C -- D -- {
                    E' -- F' -- G'
                }
            };
            % Branch
            \gitbranch
                {master}
                {above=of D}
                {D}
            % Branch
            \gitbranch
                {feature1}
                {above=of G'}
                {G'}
            % HEAD reference
            \gitHEAD
                {above=of feature1}
                {feature1}
        \end{tikzpicture}%
        }
    \end{frame}
    \begin{frame}[fragile]
        \frametitle{Git essentials}
        \framesubtitle{Cherry picking}
        \begin{lstlisting}[gobble=12]
            git cherry-pick <commit hash>
        \end{lstlisting}
        $\rightarrow$ Select commit by hash and copy it on end of current branch as new commit.\newline
        \resizebox{!}{0.25\paperheight}{%
            \begin{tikzpicture}
                % Commit DAG
                \gitDAG[grow right sep = 2em]{
                    A -- B --{
                        C -- D ,
                        E -- F -- { G }
                    } 
                };
                % HEAD reference
                \gitHEAD
                    {above=of D}
                    {D}
            \end{tikzpicture}%
        }
        \begin{center}$\Downarrow$\end{center}
        \resizebox{!}{0.25\paperheight}{%
            \begin{tikzpicture}
                % Commit DAG
                \gitDAG[grow right sep = 2em]{
                    A -- B --{
                        C -- D -- { F' },
                        E -- F -- { G }
                    } 
                };
                % HEAD reference
                \gitHEAD
                    {above=of F'}
                    {F'}
            \end{tikzpicture}%
        }
    \end{frame}

    \begin{frame}
        \frametitle{Scenarios}
        \huge{It's time to break some stuff. }
    \end{frame}
    \begin{frame}[fragile]
        \frametitle{Scenarios}
        \framesubtitle{Fresh start - cloning}
        \begin{lstlisting}[gobble=12]
            git clone https://bitbucket.honeywell.com/users/e878714/
                        repos/playground/commits
        \end{lstlisting}
        $\rightarrow$ This downloads whole repository into current working directory. (that URL should not be splitted but slide is too short...)
    \end{frame}
    \begin{frame}[fragile]
        \frametitle{Scenarios}
        \framesubtitle{Branching}
        $\rightarrow$ We do not want to fight all together so we create own branches
        \begin{lstlisting}[gobble=12]
            git branch your-branch-name
            git checkout your-branch-name
        \end{lstlisting}
        $\rightarrow$ Now we should all be at 'your-branch-name' branch.
    \end{frame}
    \begin{frame}[fragile]
        \frametitle{Scenarios}
        \framesubtitle{Committing}
        $\rightarrow$ Do some changes (anything) in code you just cloned
        \begin{lstlisting}[gobble=12]
            git commit -am "Your commit message"
        \end{lstlisting}
        $\rightarrow$ This commits all your changes with your commit message.
    \end{frame}
    \begin{frame}[fragile]
        \frametitle{Scenarios}
        \framesubtitle{Pushing}
        $\rightarrow$ Let's put our changes on server
        \begin{lstlisting}[gobble=12]
            git push
        \end{lstlisting}
        $\rightarrow$ This should actually fail because we are pushing to non-existen branch on server
        \begin{lstlisting}[gobble=12]
            git push --set-upstream origin your-branch-name
        \end{lstlisting}
        $\rightarrow$ And this will work. It creates new branch with our name on remote server named 'origin'
    \end{frame}
    \begin{frame}[fragile]
        \frametitle{Scenarios}
        \framesubtitle{Adding}
        $\rightarrow$ Let's create new file and add it to repository
        \begin{lstlisting}[gobble=12]
            git add our-new-file.txt
        \end{lstlisting}
        $\rightarrow$ This adds our-new-file.txt to staging area. We need to commit it next.
        \begin{lstlisting}[gobble=12]
            git commit -am "Your commit message"
        \end{lstlisting}
        $\rightarrow$ If we are happy with it, we can push it.
        \begin{lstlisting}[gobble=12]
            git push
        \end{lstlisting}
    \end{frame}
    \begin{frame}[fragile]
        \frametitle{Scenarios}
        \framesubtitle{Merging and pulling}
        $\rightarrow$ Now we want to merge our stuff into master. Normally we will create pull request but today we will merge direcly to master. We start with checking out master branch.
        \begin{lstlisting}[gobble=12]
            git checkout master
        \end{lstlisting}
        $\rightarrow$ Now we need to get latest changes from it.
        \begin{lstlisting}[gobble=12]
            git pull
        \end{lstlisting}
        $\rightarrow$ And now we can merge it.
        \begin{lstlisting}[gobble=12]
            git merge vojtuv-test
        \end{lstlisting}
        $\rightarrow$ And push result.
        \begin{lstlisting}[gobble=12]
            git push
        \end{lstlisting}
    \end{frame}
    \begin{frame}[fragile]
        \frametitle{Scenarios}
        \framesubtitle{Log, Checkout}
        $\rightarrow$ We want to work on our branch again (or somebody elses branch).\newline
        $\rightarrow$ We get latest changes firts from all branches and delete references to non-existent ones.
        \begin{lstlisting}[gobble=12]
            git fetch --all --prune
        \end{lstlisting}
        $\rightarrow$ Now we can take a look on current repository log
        \begin{lstlisting}[gobble=12]
            git log --graph
        \end{lstlisting}
        $\rightarrow$ And eventually checkouts something
        \begin{lstlisting}[gobble=12]
            git checkout our-branch-name
        \end{lstlisting}
    \end{frame}
    \begin{frame}[fragile]
        \frametitle{Scenarios}
        \framesubtitle{Stash}
        $\rightarrow$ Let's do some changes...\newline
        $\rightarrow$ Ok, we don't want them actaully. We need to get rid of them. But it's really nice piece of work so we will keep them for later.
        \begin{lstlisting}[gobble=12]
            git stash
        \end{lstlisting}
        $\rightarrow$ Ok, it's gone. You can verify it with status.
        \begin{lstlisting}[gobble=12]
            git status
        \end{lstlisting}
        $\rightarrow$ Now fast forward a little bit and assume we want them back. We can be on different branch, we can be in future, it doesn't matter. 
        \begin{lstlisting}[gobble=12]
            git stash list
            git stash apply stash@{0}
        \end{lstlisting}
        $\rightarrow$ Check who's back!
    \end{frame}
    \begin{frame}[fragile]
        \frametitle{Scenarios}
        \framesubtitle{Reset}
        $\rightarrow$ Oh my god, that change was rubbish... We need to delete it.
        \begin{lstlisting}[gobble=12]
            git reset --hard
        \end{lstlisting}
        $\rightarrow$ Uff, it's gone.
    \end{frame}
    \begin{frame}
        \frametitle{Scenarios}
        \framesubtitle{Git online playground}
        If you want to play with branching, rebasing, merging, squashing etc., this is great tool for playing with it.\newline
        \url{https://learngitbranching.js.org}
    \end{frame}
    \begin{frame}
        \frametitle{Troubleshooting}
        \framesubtitle{Windows Crenedital Manager}       
        Git.exe uses Credential manager for storing login information. It's nice feature with one undesired issue - password change doesn't affect stored credentials so you need to remove record manually when you change your password.
        \noindent\centerimg[width=0.6\paperwidth]{credentials.png} 
    \end{frame}

    \begin{frame}[fragile]
        \frametitle{Resources}
        \nocite{*}
        \printbibliography     
    \end{frame}
    \begin{frame}
        \Huge{Q \& A}
    \end{frame}
    \begin{frame}
        \Huge{Thank you for your attention.}
    \end{frame}
\end{document}